\documentclass[10pt,a4paper]{amsart}
\usepackage[margin=19mm]{geometry}
\usepackage{amsmath,amssymb,mathtools}
\usepackage[scr=rsfs,cal=euler]{mathalfa}
\usepackage{xfrac}
\usepackage[unicode,hidelinks,bookmarks=false]{hyperref}
\newcommand{\ie}{i.e.\ }
\newcommand{\cf}{cf.\ }
\newcommand{\resp}{respectively\ }
\newcommand{\Subsection}{\subsection}
\providecommand{\nobreakdash}{\nobreak\hbox{-}}
\setlength{\emergencystretch}{2em}
\multlinegap0pt
\usepackage{bm}
\usepackage{bbm}
\usepackage{dsfont}
\usepackage{xspace}
\usepackage{cancel}
\usepackage{mathtools}
\usepackage[shortlabels]{enumitem}
\usepackage{amsthm}
\makeatletter
\@ifundefined{theorem}{\newtheorem{theorem}{Theorem}}{}
\@ifundefined{claim}{\newtheorem{claim}[theorem]{Claim}}{}
\@ifundefined{proofclaim}{\newtheorem{proofclaim}{Proofclaim}}{}
\@ifundefined{proposition}{\newtheorem{proposition}[theorem]{Proposition}}{}
\makeatother
\makeatletter
\newcommand{\eps}{\varepsilon}
\expandafter\ifx\csname proofclaim\endcsname\relax
\newenvironment{proofclaim}[1][Proof of claim]{\begin{proof}[#1]\renewcommand*{\qedsymbol}{\(\blacksquare\)}}{\end{proof}}
\fi
\makeatother
\title[Complete tripartite subgraphs of balanced tripartite graphs ]{Complete tripartite subgraphs of balanced tripartite graphs with large minimum degree}
\author{Yihan Chen, Jialin He, Allan Lo, Cong Luo, Jie Ma, Yi Zhao}
\date{Source: arXiv:2411.19773v2 (CC BY 4.0)}
\begin{document}
\maketitle

\noindent\emph{Source and licence.} Complete tripartite subgraphs of balanced tripartite graphs with large minimum degree. Version arXiv:2411.19773v2 (CC BY 4.0), distributed under the Creative Commons Attribution 4.0 International licence, as recorded in the arXiv licence metadata for this exact version and shown on the abstract page https://arxiv.org/abs/2411.19773v2. Full rights notice: licenses/arxiv-2411.19773-CC-BY-4.0.txt.\par\smallskip

\noindent\textit{Editorial scope.} Counted unit: the Claim whose source label is \texttt{Claim:A\_i} in the article (source0.tex lines 594--618, source label \texttt{Claim:A\_i}), delivered together with the article's own complete original proof of it (source0.tex lines 620--726, one \texttt{proof} environment of 107 lines). No mathematics is written, repaired or re-derived: statement and proof are the article's own text line for line. Required context retained uncounted: Proposition: A-B set (lines 550--552); eqn:vw (lines 577--583). Every definition, display and label of those lines is retained unchanged. Omitted from this package and not redistributed: 1--549, 553--576, 584--593, 619--619, 727--1037. Nothing inside a retained excerpt is omitted or altered: every retained body is delivered exactly as the article's lines are written, so no omission receipt of any kind accompanies this package. Citations: the retained excerpts carry no citation command at all, so no bibliography entry of the article is transcribed and no reference is redistributed. Reference adaptation: none was needed and none was made; every reference inside the delivered unit resolves inside it, so no unresolved reference or citation is produced. Printed slips kept literal: no printed word, symbol, hypothesis or number of the counted statement or of its proof was altered. Layout and class: the article's own class is \texttt{article} with packages including enumerate, amssymb, a4wide, latexsym, makeidx, epsfig, amsthm, dsfont, amsmath, lipsum, mathtools, enumitem, mathrsfs, xcolor; the wrapper is the lane's pinned \texttt{amsart} editorial wrapper at 10pt on A4, which restates the theorem-like environments the retained text needs and adds the article's own macro definitions that the retained text needs (\texttt{\textbackslash eps}), with extra wrapper packages (bm, bbm, dsfont, xspace, cancel, mathtools, enumitem). The delivered numbering is the unit's own. Assets: no retained span contains a figure environment, a graphics-inclusion command, a PDF-page inclusion, a table or a TikZ picture; no image file is copied into this package, the package declares no asset dependency and no image file is redistributed with it. Licence: version arXiv:2411.19773v2 of this article is distributed under the Creative Commons Attribution 4.0 International licence, as recorded in the arXiv licence metadata for this exact version and shown on the abstract page https://arxiv.org/abs/2411.19773v2. A full rights notice is delivered at licenses/arxiv-2411.19773-CC-BY-4.0.txt. Characters: every retained excerpt is pure ASCII, so the delivered engine, XeLaTeX with its default Latin Modern text font, reports no missing character.
	\begin{proposition}\label{Proposition: A-B set}
		Let $0\leq \lambda\leq 1/{10}$ and $G$ be a bipartite graph with vertex classes $A$ and~$B$. Suppose that for all $a\in A$, $d_G(a)\geq (1-\lambda)|B|$. Then there exists a subset $B'\subseteq B$ of size $|B'|\geq (1-5\lambda)|B|$ such that for all $a\in A$ and $b\in B'$, $d_G(a,B')\geq (1-2\lambda)|B'|$ and $d_G(b,A)\geq 4|A|/5$.
	\end{proposition}

\begin{align}
\alpha n &\ge T(vw) = | N(w)\cap N(v) | \ge d^+(w) + d^-(v) - n 
= d(w) - d^-(w) + d^-(v) - n \nonumber \\
& \ge \delta(G) - d^-(w) + d^-(v) - n 
\ge d^-(v) - d^-(w), \nonumber \\
			d^-(w)& \geq d^-(v)-\alpha n.	\label{eqn:vw}
	\end{align}

		\begin{claim}\label{Claim:A_i}
			The sets $A_i$'s satisfy the following properties.
			\begin{enumerate}[label = {\rm (\roman*)}]
				\item For $i\in [12]$, we have $|A_i|\geq \beta n$.
				\label{itm:C_6:1}
				
				\item For $i\in \{0,1,\dots,12\}$ and $v\in A_i$, we have $d^-(v)\geq (1-\beta-i\alpha)n$.
				\label{itm:C_6:2}
				
				\item For $i\in \{0,\dots,11\}$ and $v\in A_i$, we have $d(v, A_{i+1})\ge |A_{i+1}|-(i+1)\alpha n\geq (\beta -(i+1)\alpha)n$ and $d^+(v)\leq (\beta+i\alpha)n$. In particular, for $i\in [12]$, $|A_i|\leq d^+(a_{i-1})\leq (\beta +(i-1) \alpha)n$.
				\label{itm:C_6:3}
				
				\item For $i\in \{0,1,\dots,10\}$ and $v\in A_i$, we have $d^-(v,A_{i+2})\leq (i+3)\alpha n$.
				\label{itm:C_6:4}
				
				\item For $i\in \{0,1,\dots,9\}$, $|A_i\cap A_{i+3}|\leq 5\sqrt{\alpha} n$.
				\label{itm:C_6:5}
				
				\item For $i\in \{0,1,\dots,10\}$ and $v\in A_i$, we have $|V_{i-1}\setminus\left(A_{i+2}\cup N^-(v)\right)|\leq (2i+3)\alpha n$.
				\label{itm:C_6:6}
				
				\item For $i\in \{1,\dots,6\}$, $|A_i\cap A_{i+6}|\geq |A_i|-8\sqrt{\alpha} n\geq (\beta -8\sqrt{\alpha})n$.
				\label{itm:C_6:7}
			\end{enumerate}
		\end{claim}

		\begin{proofclaim}
			We have $|A_i| = d^+(a_{i-1})\ge \delta^+(G) = \beta n$ giving \ref{itm:C_6:1}.
			
			We prove \ref{itm:C_6:2} by induction on~$i$.
			For $i=0$, by using $\beta n = d^+(a_0)$, we have
			\begin{align*}
				d^-(a_0) = d(a_0) - d^+(a_0) \geq n - \beta n  = (1-\beta)n.
			\end{align*}
			So we may assume $i \in [12]$.
			For $v\in A_i$, note that $a_{i-1}v \in E(G)$.
			Together with~\eqref{eqn:vw}, we have 
			\begin{align*}
				d^-(v)\geq d^-(a_{i-1})-\alpha n\geq (1-\beta- (i-1)\alpha )n - \alpha n = (1-\beta-i\alpha) n.
			\end{align*}
			Hence \ref{itm:C_6:2} holds.
			
			For $i\in \{0,\dots,11\}$ and $v\in A_i$, we first show $d^+(v)\leq (\beta+i\alpha)n$ by using \ref{itm:C_6:2}. We know $d^+(a_0)= \beta n$. For $i\in [11]$, since $v\in A_i = N^+(a_{i-1})$, we have $|N^+(v)\cap N^-(a_{i-1})| = T( v a_{i-1}) \le \alpha n$. Consequently,
			\begin{align*}
				d^+(v) \le |N^+(v)\cap N^-(a_{i-1})| + |V_{i+1}\setminus N^-(a_{i-1})| 
				\overset{\mathclap{\text{\ref{itm:C_6:2}}}}{\le}
				\alpha n + (\beta + (i-1) \alpha)n = (\beta+i \alpha)n.
			\end{align*}
			This implies that $|A_i|\leq d^+(a_{i-1})\leq (\beta+(i-1)\alpha)n$ for $i\in [12]$.

			Next we show $d(v, A_{i+1})\ge |A_{i+1}|-(i+1)\alpha n\geq (\beta -(i+1)\alpha)n$ for $v\in A_i$ and $i\in \{0,\dots,11\}$. First, $d(a_0,A_1)=|A_1| > (\beta - \alpha) n$. 
			Note that 	
			\begin{align}
				\nonumber
				| V_{i+1}\setminus  ( N^-(a_{i-1}) \cup  N^+(v) )|
				& = |V_{i+1}\setminus N^-(a_{i-1})|  - |N^+(v) \setminus N^-(a_{i-1}))|\\
				\nonumber
				& \overset{\mathclap{\text{\ref{itm:C_6:2}}}}{\le }
				(\beta+ (i-1) \alpha)n - d^+(v) + T(a_{i-1} v)\\
				& \le (\beta+ (i-1) \alpha)n - \delta^+(G) + \alpha n 
				= i \alpha n.
				\label{eqn:C_6:1}
			\end{align}
It follows that 
			\begin{align*}
				d(v,A_{i+1})  
				& \ge |A_{i+1}| - |A_{i+1}\cap N^-(a_{i-1})| - |A_{i+1} \setminus  ( N^-(a_{i-1}) \cup  N^+(v) )|\\
				& \ge |A_{i+1}| - T(a_{i}a_{i-1}) - |V_{i+1} \setminus  ( N^-(a_{i-1}) \cup  N^+(v) )|\\
				& \overset{\mathclap{\text{\eqref{eqn:C_6:1}}}}{\ge }
				|A_{i+1}| - \alpha n - i \alpha n \overset{\mathclap{\text{\ref{itm:C_6:1}}}}{\ge } (\beta - (i+1)\alpha ) n.
			\end{align*}
			Hence \ref{itm:C_6:3} holds.
	
		Suppose $i\in \{0,1,\dots,10\}$ and $v\in A_i$. Let $w\in N^+(v,A_{i+1})$, which exists because \ref{itm:C_6:3} implies that $d^+(v, A_{i+1})\ge |A_{i+1}| - (i+1)\alpha n >0$. Then 
			\begin{align*}
				\alpha n \ge T(vw) &= |N^-(v)\cap N^+(w)|
				\geq |N^-(v)\cap N^+(w)\cap A_{i+2}| \\
				&\geq d^-(v,A_{i+2}) + d^+(w,A_{i+2})-|A_{i+2}| \\
				& \overset{\mathclap{\text{\ref{itm:C_6:3}}}}{ \ge } d^-(v,A_{i+2}) +(|A_{i+2}|-(i+2)\alpha n)-|A_{i+2}| = d^-(v,A_{i+2}) -(i+2)\alpha n.
			\end{align*}			
			This gives $d^-(v,A_{i+2})\le (i+3)\alpha n$, confirming 
			\ref{itm:C_6:4}. 			
			
			By~\ref{itm:C_6:1} and~\ref{itm:C_6:3}, for all $v\in A_{i+2}$, we have 
			\begin{align*}
				d^+(v,A_{i+3})>|A_{i+3}|-(i+3)\alpha n \geq (1-\sqrt{\alpha})|A_{i+3}|,
			\end{align*} 
			here we use the assumption that $\alpha = (\eps/ 30)^2$ and $\beta \ge 30\sqrt{\alpha}$. 
			By applying Proposition~\ref{Proposition: A-B set} on $G[A_{i+2},A_{i+3}]$, we have a subset $A_{i+3}'\subseteq A_{i+3}$ with size $(1-5\sqrt{\alpha})|A_{i+3}|$ such that, for all~$w\in A_{i+3}'$, 
			\begin{align*}
				d(w,A_{i+2})\geq \frac{4}{5}|A_{i+2}| > (i+3)\alpha n.
			\end{align*}
			Since $d(v, A_{i+2})\le (i+3)\alpha n$ for all $v\in A_i$ by~\ref{itm:C_6:4}, it follows that $A_i\cap A_{i+3}\subseteq A_{i+3}\setminus A_{i+3}'$. Therefore, $|A_i\cap A_{i+3}| \le 5\sqrt{\alpha} |A_{i+3}|\le 5\sqrt{\alpha} n$ confirming \ref{itm:C_6:5}.
			
			For all $v\in A_i$,
			\begin{align*}
				d^-(v,V_{i-1}\setminus A_{i+2})& = d^-(v)-d^-(v,A_{i+2})
				\overset{\mathclap{\text{\ref{itm:C_6:2},\ref{itm:C_6:4}}}}{ \ge }
				(1-\beta-i\alpha )n-(i+3) \alpha n\\
				&= (1-\beta-(2i+3) \alpha )n
				\overset{\mathclap{\text{\ref{itm:C_6:1}}}}{ \ge }
				|V_{i-1}|-|A_{i+2}|-(2i+3)\alpha n\\
				& = |V_{i-1}\setminus A_{i+2}|-(2i+3)\alpha n.
			\end{align*}
			Hence $|V_{i-1}\setminus\left(A_{i+2}\cup N^-(v)\right)|\leq (2i+3)\alpha n$ and \ref{itm:C_6:6} holds.
			
			Finally, for $0\le i\le 6$ and $v\in A_i$, we have
			\begin{align*}
				d^-(v,A_{i+5})& \geq  d^-(v,A_{i+5}\setminus A_{i+2})
				\geq |A_{i+5}\setminus A_{i+2}|-|V_{i-1}\setminus\left(A_{i+2}\cup N^-(v)\right)|\\
				&			\overset{\mathclap{\text{\ref{itm:C_6:5},\ref{itm:C_6:6}}}}{ \ge }
				|A_{i+5}|-5\sqrt{\alpha} n-(2i+3)\alpha n
				\geq  |A_{i+5}|-6\sqrt{\alpha } n.
			\end{align*}
			Hence there exists a vertex $w\in A_{i+5}$ such that
			\begin{align*}
				d^+(w,A_i)&\geq \frac{e(A_{i},A_{i+5})}{|A_{i+5}|} \ge 		
				\frac{|A_i|(|A_{i+5}|-6\sqrt{\alpha} n)}{|A_{i+5}|}
				=|A_i|-\frac{|A_i|}{|A_{i+5}|}6\sqrt{\alpha} n
				\geq |A_i|-7\sqrt{\alpha} n,
			\end{align*}
			where the last inequality holds because $|A_{i+5}|\geq \beta n$ and $|A_i|\leq (\beta + (i-1)\alpha)n$ by~\ref{itm:C_6:1} and~\ref{itm:C_6:3}, and consequently, $\frac{|A_i|}{|A_{i+1}|}\le \frac{\beta + 5\alpha}{\beta}<\frac76$ by our assumption on $\alpha$ and $\beta$. 
			Therefore,
			\begin{align*}
				|A_i\cap A_{i+6}|&\geq |A_i\cap A_{i+6}\cap N^+(w)|
				\geq d^+(w,A_i)+d^+(w,A_{i+6})-d^+(w)\\
				&
				\overset{\mathclap{\text{\ref{itm:C_6:3}}}}{>}
				|A_i|-7\sqrt{\alpha} n + (\beta -(i+6)\alpha)n -(\beta+(i+5)\alpha)n
				\geq |A_i|-8\sqrt{\alpha} n,
			\end{align*}
			confirming \ref{itm:C_6:7}.
		\end{proofclaim}

\end{document}
